\begin{abstract}
经典 Tucker 分解以多线性低秩结构为核心，但当观测张量来源于非线性生成、饱和响应、复杂高阶交互或异方差扰动时，纯多线性假设往往会产生系统性失配。针对这一问题，本文提出多项式链接的非线性 Tucker 分解模型 NTD-PL（Nonlinear Tucker Decomposition with Polynomial Links）。该模型保留共享的 Tucker 潜在张量作为低秩骨架，并在其上施加显式的多项式链接，从而将标准 Tucker 作为线性特例纳入统一框架，并以少量链接系数实现对高阶交互的受控扩展。

理论上，本文说明了 NTD-PL 与 Tucker 的嵌套关系，给出了多项式生成情形下的精确表达结果，以及解析函数情形下的有限阶截断逼近界；同时构建了统一的带掩码学习目标，并结合 Tucker 初始化、梯度更新与多项式系数的更新，给出了可实现的交替优化方法。实验结果表明：在线性源数据下，NTD-PL 与 Tucker 基本重合，不会制造伪增益；在统一非线性残差协议下，NTD-PL 在结构匹配与有限阶逼近两类情形中均取得稳定收益；在真实高光谱重构与随机缺失补全任务上，NTD-PL 在同参数预算下持续优于 Tucker。进一步的机制对照表明，这种收益既不能归因于单纯增参，也不能由对 Tucker 输出施加一次共享多项式校正所替代。

总体而言，本文为保持 Tucker 骨架而增强非线性生成机制提供了一种显式、低维、可分析的结构化方案。
\end{abstract}
